% Root File for Prentice-Hall Book
%         11 point version
% Root file for Pearson Generic book
%   **** Trim size 9 x 7 ****
%
%   Prepared for LaTeX2e February, 1997
%   Derived from LaTeX 2.09 version
%  ** Updated for Style June 2002 **
%  ** Updated for Style January 2010 **
%  ** Updated with minor corrections Nov 2015 **

%%%%%%%%%% First, determine and select the trim size of the book
%
% Pearson uses three basic trim sizes.  Ask your Production Editor which one will be used for your book.
% If your book will be 7 inches wide, use that documentclass statement and comment out the other one.
% If your book will be 7 and 3/8 inches wide, use that documentclass statement and
%            comment out the other one.

\documentclass[twoside,10pt,letterpaper,usenames]{newstyle-PearsonGeneric-7-38}
\usepackage[twoside]{geometry} % to set the dimensions of the page

%%%%%%%% Then, also select the trim size of the book here.
%
%    Uncomment the appropriate trim size for your book, and comment-out the others.
%%%%%%%%%%%%%% Trim Size 3:  7 3/8 x 9 1/8
\geometry{                                  % setting all the necessary dimensions...
paperwidth=7.375in,                  % Do not change these settings.
paperheight=9.125in,                 % These set up the trim size for the book.
lmargin=1in,
rmargin=.875in,
bmargin=.650in,
tmargin=.95in,
width=5.5in,
height=7.4in,
marginparwidth=0in,
marginparsep=0in,
headheight=0.2in,
headsep=.25in,
footskip=.025in}
%%%%%%%%%%%%%%%%%%%%%

%%%%%%%% Skip over this part.  These lines should remain unaltered.
%

\usepackage{amsmath} % to ensure attractive mathematics
\usepackage{amsfonts} % to ensure attractive mathematics
\usepackage{amssymb} % to ensure attractive mathematics

%JMB: Using Cousine as code font
\usepackage{mathspec}
\setmonofont[Ligatures=NoCommon,Scale=.9,SlantedFont={Cousine-Italic},BoldFont={Cousine Bold},BoldSlantedFont={Cousine-Bold Italic}]{Cousine-Regular} % to set our code font
\newfontface{\SecCode}[Ligatures=NoCommon,Scale=1.3]{Cousine-Bold}
\newfontface{\SubsecCode}[Ligatures=NoCommon,Scale=1.1]{Cousine-Bold}
\newfontface{\SubsubsecCode}[Ligatures=NoCommon,Scale=.91]{Cousine-Bold}
\newfontface{\ParaCode}[Ligatures=NoCommon,Scale=.9]{Cousine-Bold}
\newfontface{\SubparaCode}[Ligatures=NoCommon,Scale=.9]{Cousine-Regular}

%%    These lines input packages that are used for this template.
\usepackage[letter,cam,center]{crop-AlteredCamMarks}  % to add trim marks; updated Nov 2015 to allow more space between marks and trim edge
\usepackage{graphicx} % for graphics layout
\usepackage{float}      % to allow things to float even if they normally don't
\usepackage{wrapfig} % to allow text to wrap around small figures
\usepackage{outline} % to allow outline creation
\usepackage{framed-PearsonGeneric} % to create the leftside vertical bar for the example environment
% The listings package isn't included in the template folder because
%       it is part of any up-to-date TeX installation
\usepackage[writefile]{listings}  % for setting attractive computer code
\lstset{basicstyle=\ttfamily\small,framerule=0.5pt} % setting the code in small sans serif font and
                                                    % setting the margin for the code.  Nov 2015: setting the rules to be 0.5pt
                                                    % JMB: Changed font to Cousine, removed margin for Lakos books
\usepackage[stable,bottom]{footmisc} % to stabilize the footnote environment and to ensure footnotes
                                     % appear at the bottom of the page, even if a figure or table
                                     % appears at the bottom
\usepackage[flushleft]{threeparttable} % to set attractive footnotes within table environments
\usepackage{titleref} % to allow cross-referencing by title
\usepackage{xr} % to allow cross-referencing between files
\usepackage{colortbl} % to allow the use of color/shading in tables
\usepackage{caption} % to set the appearance of captions
\captionsetup[table]{singlelinecheck=off,labelfont={sf,bf,small},font={sf,bf,small}}
\captionsetup[figure]{labelfont={sf,bf,small},font={sf,bf,small}}
\captionsetup[lstlisting]{labelfont={sf,bf,small},font={sf,bf,small}}
\usepackage{tocloft-hacked-PearsonGeneric} % to allow manipulation of the TOC


%%%%%%%% Determine your indexing plan.
%
% Most of the time, the publisher will hire an independent indexing expert to create an
%    index for your book.  Your production editor can let you know the plan for your book.
% Use the line below only if your production plan includes the creation of a LaTeX index.
% Note that this situation is uncommon.
\makeindex           % if your book will include a LaTeX-generated index
\usepackage{See-makeidx}

%%%%%%%%%% These lines set up the Preamble of the document. You can skip over this part.
%% Preamble:
% New commands and/or command redefinitions
%
% You can also place such commands in a macro file (e.g. mydefs.tex)
% and load them in the preamble with "input mydefs"
% Here are some examples:

\newcommand{\be}{\begin{equation}}           %---- begin numbered equation
\newcommand{\ee}{\end{equation}}             %---- end numbered equation
\newcommand{\dst}{\everymath{\displaystyle}} %---- use displaystle in eqs.
\newcommand{\Nuclide}[2]{${}^{#2}$#1}        %---- Nuclide macro
\renewcommand{\SS}[1]{${}^{#1}$}             %---- superscript in text
\newcommand{\STRUT}{\rule{0in}{3ex}}         %--- small strut
\newcommand{\dotprod}{{\scriptscriptstyle \stackrel{\bullet}{{}}}}
\renewcommand{\tilde}{\widetilde}
\renewcommand{\hat}{\widehat}
\renewcommand{\theequation}{\arabic{chapter}.\arabic{equation}}
\renewcommand{\thefigure}{\arabic{chapter}--\arabic{figure}}
\renewcommand{\thetable}{\arabic{chapter}--\arabic{table}}
\newcommand{\ul}{\underline}
\newcommand{\dul}{\underline{\underline}}
\newcommand{\ol}{\overline}
\newcommand{\beq}{\begin{equation}}
\newcommand{\eeq}{\end{equation}}
\newcommand{\bea}{\begin{eqnarray}}
\newcommand{\eea}{\end{eqnarray}}
\newcommand{\beas}{\begin{eqnarray*}}
\newcommand{\eeas}{\end{eqnarray*}}
\newcommand{\ba}{\begin{array}}
\newcommand{\ea}{\end{array}}

%%%%%%%% set chapter author commands
\newcommand{\chapauthor}[1]{\hspace*{\fill}{\textsf{\large 
 #1}}\\[1.5ex]}
%%%%%%%%%%%%%%

%%%%%%% set epigraph commands
\newcommand{\epigraph}[1]{\vspace*{2ex}{\raggedleft{\textit{#1}}\\[6ex]}}
%%%%%%%%%%%%

%%%%%% set definition commands
\newcounter{definitionctr}[chapter]
\def\thedefinitionctr{\arabic{chapter}--\arabic{definitionctr}}
\newcommand{\definition}[1]{\refstepcounter{definitionctr}\par\vskip2ex \indent\parbox[]{.96\textwidth}{\indent{\textbf{Definition~\thedefinitionctr.}~~#1}\par\vskip2ex}}
%%%%%%%%%%%%%%%%%

%%%%%%%%%%%%%%% create a thick line for tables
\newcommand{\thickhline}{\noalign{\hrule height 1.5pt}} 
%%%%%%%%%%%%%

%%%%%%%%%%%%%% define the unnumbered list
\newenvironment{unnumlist}{\begin{list}{}% empty for no label
{}} % empty since we don't need a counter and the standard lengths/indentations are fine
{\end{list}} 
%%%%%%%%%%%%%%%%%%%

%%%%%%%%%%% define a simple command for the quote environment
\newenvironment{italquote}{\begin{quote}{}% empty for no label
{}\itshape} % empty since we don't need a counter and the standard lengths/indentations are fine
{\end{quote}} 
%%%%%%%%%%%

%%%%%%%%%%% define a simple command for the quote environment
\newenvironment{dialogue}[2]{\begin{quote}{\bfseries #1:~}% empty for no label
{} #2} % empty since we don't need a counter & the standard lengths/indentations and text style are fine
{\end{quote}} 
%%%%%%%%%%%

%%%%%%%%%%%%% define the multicolumn list environment
\newcommand{\mclhead}{\bfseries} % define a multi-column list (mcl) heading
\newlength{\mcltopsep}% define a spacing command
\setlength{\mcltopsep}{\topsep} % 
\newenvironment{mcl}{\par\vspace*{\mcltopsep}
\begin{tabular}{ll}}{\end{tabular}\linebreak[4] \par\vspace*{.5\mcltopsep}}
%%%%%%%%%%%%%%%%

%%%%%%%%%%%% define some sidebar commands
% the sidebar is created via the framed.sty package
\definecolor{shadecolor}{gray}{0.9} % define the color to be used for the sidebar
% define the sidebar head  ---  corrected bug 4/19/2011
\newcommand{\sidebarhead}[1]{\noindent{\cmssbxsection #1}\hspace*{\fill}\linebreak\nopagebreak} 
%%%%%%%%%%%%%%

%%%%%%%%%%% define note, tip, and warning environments
\newcommand{\note}[2]{\pagebreak[3]\par\vspace*{2ex} % define a new note command 
\noindent{\sffamily\bfseries\small NOTE}\linebreak\vspace*{-1.5ex}\nopagebreak
\rule[1.5ex]{\textwidth}{.75pt}\linebreak\nopagebreak
{\sffamily\small  {\bfseries{\raggedright #1}\\}\nopagebreak
\noindent{#2}}\par
\noindent\rule[1ex]{\textwidth}{.75pt}\par\vspace*{1ex}\pagebreak[3] 
}

\newcommand{\tip}[2]{\pagebreak[3]\par\vspace*{2ex} % define a new tip command 
\noindent{\sffamily\bfseries\small TIP}\linebreak\vspace*{-1.5ex}\nopagebreak
\rule[1.5ex]{\textwidth}{.75pt}\linebreak\nopagebreak
{\sffamily\small  {\bfseries{\raggedright #1}\\}\nopagebreak
\nopagebreak
\noindent{#2}}\par
\noindent\rule[1ex]{\textwidth}{.75pt}\par\vspace*{1ex}\pagebreak[3]
}

\newcommand{\warning}[2]{\pagebreak[3]\par\vspace*{2ex} % define a new warning command 
\noindent{\sffamily\bfseries\small WARNING}\linebreak\vspace*{-1.5ex}\nopagebreak
\rule[1.5ex]{\textwidth}{.75pt}\linebreak\nopagebreak
{\sffamily\small  {\bfseries{\raggedright #1}\\}\nopagebreak
\noindent{#2}}\par
\noindent\rule[1ex]{\textwidth}{.75pt}\par\vspace*{1ex}\pagebreak[3]
}
%%%%%%%%%%%%%%

%%%%%%%%%%% define glossary command
\newcommand{\gloss}[2]{\pagebreak[3]\par
\noindent{\sffamily\bfseries #1}\\
\noindent{#2}\hspace*{\fill}\\ 
\par\pagebreak[3]}
%%%%%%%%%%%%%%%%
  % Loading the definitions needed for this book

%
% This file is included before the document is begun.  All entries put in here
% likely need to be approved by the publisher.
%

\newcommand{\todo}[1]{}

% C++ deserves being well-typeset in a fashion matching the standard document
\newcommand{\Cpp}{C\kern-0.2em\protect\raisebox{.15ex}{+\kern-0.2em+}}
\newcommand{\CppIII}{\Cpp{} 2003}
\newcommand{\CppXI}{\Cpp{} 2011}
\newcommand{\CppXIV}{\Cpp{} 2014}
\newcommand{\CppXVII}{\Cpp{} 2017}
\newcommand{\CppXX}{\Cpp{} 2020}

\definecolor{maroon}{cmyk}{.26,1,.63,.17}
\definecolor{skyblue}{cmyk}{.62,.07,0,0}
\definecolor{forestgreen}{cmyk}{.9,0,.9,0}
\definecolor{mustard}{cmyk}{.22,.14,1,0}
\definecolor{purple}{cmyk}{.83,.81,0,0}
\definecolor{codegray}{cmyk}{0,0,0,.70}

%%% PRINT BOOK SETTINGS
%\lstset{language=C++,basicstyle=\ttfamily\small,framerule=0.5pt,
%emph=[1]{const},
%  emphstyle=[1]{\color{forestgreen}\ttfamily},
%  moreemph=[2]{int,void,char},
%  emphstyle=[2]{\color{forestgreen}\ttfamily},
%    moreemph=[3]{decltype},
%      emphstyle=[3]{\color{forestgreen}\ttfamily},
%    keywordstyle={\color{forestgreen}\ttfamily\bfseries},
%  stringstyle={\color{forestgreen}},
%commentstyle={\color{codegray}\itshape},
%numberstyle=\color{forestgreen},
%xleftmargin=2ex,
%    escapeinside={(ù}{ù)},
%    showstringspaces=false, 
%} % setting the code in small sans serif font and
% setting the margin for the code.  Nov 2015: setting the rules to be 0.5pt
%%%%%%%%%%%%%
%\newcommand{\codeincomments}{\color{codegray}\ttfamily}
%%%%%%%%%%%%%%%%%%%%%%%%%%%%%


%%% EBOOK SETTINGS
\lstset{language=C++,basicstyle=\ttfamily\small,framerule=0.5pt,
emph=[1]{const},
  emphstyle=[1]{\color{mustard}\ttfamily},
  moreemph=[2]{int,void,char},
  emphstyle=[2]{\color{maroon}\ttfamily},
    moreemph=[3]{decltype},
      emphstyle=[3]{\color{forestgreen}\ttfamily},
    keywordstyle={\color{forestgreen}\ttfamily\bfseries},
  stringstyle={\color{purple}},
commentstyle={\color{skyblue}\itshape},
numberstyle=\color{purple},
xleftmargin=2ex,
    escapeinside={(ù}{ù)},
    showstringspaces=false, 
} % setting the code in small sans serif font and
%% setting the margin for the code.  Nov 2015: setting the rules to be 0.5pt
%%%%%%%%%%%%%%
\newcommand{\codeincomments}{\color{skyblue}\ttfamily}

% All code samples use this block to mark for other build modes that they
% should be written out for compiling
%

\lstnewenvironment{cppcodeblock}[1]
{\lstset{language=[11]C++}%
}
{
}

\newcommand{\cppcode}[1]{\lstinline[language=C++]{#1}}


%%%%%%%%% Now the document begins

\begin{document}
%%%%%% commands to fine-tune to table of contents
%\cftnodots
%\renewcommand{\@dotsep}{10000}


\renewcommand{\cfttoctitlefont}{\cmsschapname}
\renewcommand{\cftpartfont}{\cmssbxparttoc}
\renewcommand{\cftchappagefont}{\cmssbxparttoc}
\renewcommand{\cftchapfont}{\cmssbxchaptoc}
\renewcommand{\cftchapnumwidth}{5.5em}
\renewcommand{\cftchappagefont}{\cmssbxchaptoc}
\renewcommand{\cftsecfont}{\sf}
\renewcommand{\cftsecindent}{0em}
%\renewcommand{\cftsecnumwidth}{0em}
\renewcommand{\cftsecpagefont}{\sf}
\renewcommand{\cftsubsecfont}{\sf}
\renewcommand{\cftsubsecindent}{4em}
%\renewcommand{\cftsubsecnumwidth}{0em}
\renewcommand{\cftsubsecpagefont}{\sf}
\renewcommand{\contentsname}{Contents}


%\setcounter{secnumdepth}{0} 
%\renewcommand{\numberline}[2]{} 
%%%%% end toc commands % Loading a few commands related to the tocloft package

%%%%%%%%%  You will need to enter the file names of all your chapter files here.
% These lines tell the xr package what chapters need to cross-reference one another.
% You'll need to add your chapter file names here
\externaldocument{cawp-application-developers}
\externaldocument{cawp-library-writers}
\externaldocument{cawp-writing-allocators}
\externaldocument{cawp-advanced}
\externaldocument{cawp-appendix}
\externaldocument{common-peerson-preamble}

%%%%% The front matter begins:
\pagenumbering{roman}             % Roman numbering
\setcounter{page}{1}              % Starting page number (specified by the publisher)
\cleardoublepage                  % This command is defined in our documentclass file. It tells
                                  % our compiler that we want the next item
                                  % to start on a right-hand page, and it inserts a blank left-hand
                                  % page if needed.
\thispagestyle{empty}
\begin{center}
{\sfititle \Cpp{} Allocators for the Working Programmer}
\end{center}
              % The half-title page
../../src/books/cawp/cawp-series.tex                  % A placeholder for the series page. The publisher normally
                                  % creates this page and just adds it into the final PDF.
\thispagestyle{empty}
\begin{center}
{\sfititle \Cpp{} Allocators for the Working Programmer}\\
\todo{Identify if we do want a subtitle}
\vspace*{54pt}
%{\sfihalftitle Subtitle of the Book}\\
\vspace*{\fill}    % This space will automatically adjust
\noindent{\sfiauthor John Lakos}\\[3ex]
\noindent{\sfiauthor Joshua Berne}\\        % Add or delete lines as necessary for author names

\vspace*{\fill}    % This space will automatically adjust
% Authors:  For the next line, use either the AW line or the PH line per your 
%   production editor's instructions. You will find both PC and Mac versions of the logos 
%   in the template folder. Use the one appropriate for your system.
\todo{PH or AW?}
%\includegraphics{PH_title_page_Mac} \\   %PH Mac
%\includegraphics{PH_stack_BLACK-PC} \\  %PH PC
\includegraphics{AWlogoTitlePage-Mac}\\   % AW Mac
%\includegraphics{AWlogoTitlePage_PC}  \\  %AW PC
{\footnotesize 
Boston~$\bullet$~Columbus~$\bullet$~Indianapolis~$\bullet$~New York~$\bullet$~San Francisco~$\bullet$~Amsterdam~$\bullet$~Cape Town\\[1ex]
Dubai~$\bullet$~London~$\bullet$~Madrid~$\bullet$~Milan~$\bullet$~Munich~$\bullet$~Paris~$\bullet$~Montreal~$\bullet$~Toronto~$\bullet$~Delhi~$\bullet$~Mexico City\\[1ex]
Sao Paulo~$\bullet$~Sidney~$\bullet$~Hong Kong~$\bullet$~Seoul~$\bullet$~Singapore~$\bullet$~Taipei~$\bullet$~Tokyo\\[1ex]
}
\end{center}
                   % The title page
\thispagestyle{empty}
\vspace*{\fill}
{\footnotesize 
\noindent Many of the designations used by manufacturers and sellers to distinguish their products are claimed as trademarks. Where those designations appear in this book, and the publisher was aware of a trademark claim, the designations have been printed with initial capital letters or in all capitals.\\
\hspace{\fill}\\
\noindent The authors and publisher have taken care in the preparation of this book, but make no expressed or implied warranty of any kind and assume no responsibility for errors or omissions. No liability is assumed for incidental or consequential damages in connection with or arising out of the use of the information or programs contained herein.\\   % change to author (singular) if this book has a single author
\hspace{\fill}\\
\noindent For information about buying this title in bulk quantities, or for special sales opportunities (which may include electronic versions; custom cover designs; and content particular to your business, training goals, marketing focus, or branding interests), please contact our corporate sales department at \mbox{corpsales}\linebreak[3]@\linebreak[3]pearsoned\linebreak[3].com or (800) 382-3419.\\
\hspace{\fill}\\
\noindent For government sales inquiries, please contact governmentsales@pearsoned.com.\\
\hspace{\fill}\\
\noindent For questions about sales outside the United States, please contact international@pearsoned.com.\\
\hspace{\fill}\\
\noindent Visit us on the Web: informit.com/aw\\     % use this line if you are publishing with AW
%\noindent Visit us on the Web: informit.com/ph\\   % use this line if you are publishing with PH
\hspace{\fill}\\
\noindent\textit{Library of Congress Cataloging-in-Publication Data}\\[.75ex]  % update this with info provided by your production editor
\textbf{LIBRARY OF CONGRESS CIP DATA WILL GO HERE; MUST BE ALIGNED AS INDICATED BY LOC}\\[.75ex]
\noindent Copyright~\copyright~2016 Pearson Education, Inc.\\
\hspace{\fill}\\
\noindent All rights reserved. Printed in the United States of America. This publication is protected by copyright, and permission must be obtained from the publisher prior to any prohibited reproduction, storage in a retrieval system, or transmission in any form or by any means, electronic, mechanical, photocopying, recording, or likewise. For information regarding permissions, request forms and the appropriate contacts within the Pearson Education Global Rights \& Permissions Department, please visit www.pearsoned.com/permissions/.\\[.75ex]

\noindent ISBN-13: NUMBER HERE\\[-.5ex] % update this with info provided by your production editor 
\noindent ISBN-10: NUMBER HERE \\[-.5ex] % update this with info provided by your production editor
\noindent Text printed in the United States on recycled paper at PRINTER INFO HERE.\\[-.5ex] % update this with info provided by your production editor
\noindent First printing, MONTH YEAR % update this with info provided by your production editor
}
                     % The cataloging-in-publication page
\thispagestyle{empty}
\vspace*{\fill}
\begin{center}This is John's dedication to Josh for being so great and \\
writing this book so well.\\
JL\\
\vspace*{3pc}
This is Josh's dedication to his wife, child, \\
and mother-in-law for being all supportive and \\
wonderful.  And to steak.  Steak is great. \\[1ex]
JMB
\end{center}
\vspace*{\fill}
              % The dedication page
\cleardoublepage
\parskip .5ex                     % Add space between Contents items
\tableofcontents                  % Make the Table of Contents
\cleardoublepage
\parskip 0in                      % Reset to zero interpar spacing

\cleardoublepage
\chapter*{Foreword}

The text of the foreword will go here.                  % Add foreword
\cleardoublepage
\chapter*{Preface}

The text of the preface will go here.  

                 % Add preface
../../src/books/cawp/cawp-ack.tex                     % Add acknowledgements
\cleardoublepage
\chapter*{About the Authors}

 % wrapfig doesn't like to be immediately after a chapter/section command 
 % so we are tricking it a bit with the \hspace and the \par commands.
\hspace*{\fill}\\  
\par\begin{wrapfigure}[12]{L}[0pt]{0pt}\includegraphics{cawp-authorphoto-JL}%
\end{wrapfigure}\noindent John Lakos, author of Large-Scale C++ Software Design
[Pearson, 1996] and Large-Scale C++ — Volume I: Process and Architecture
[Pearson, 2019],    serves at Bloomberg in New York City as a senior architect
and mentor for C++ software development worldwide. He is also an active voting
member of the C++ Standards Committee’s Evolution Working Group. From 1997 to
2001, Dr. Lakos directed the design and development of infrastructure libraries
for proprietary analytic financial applications at Bear Stearns. From 1983 to
1997, Dr. Lakos was employed at Mentor Graphics, where he developed large
frameworks and advanced ICCAD applications for which he holds multiple software
patents. His academic credentials include a Ph.D. in Computer Science (1997)
and an Sc.D. in Electrical Engineering (1989) from Columbia University. Dr.
Lakos received his undergraduate degrees from MIT in Mathematics (1982) and
Computer Science (1981).


\vspace*{18pt} % we add some space for a nice layout
\begin{wrapfigure}[12]{L}[0pt]{0pt}\includegraphics{cawp-authorphoto-JMB}%
\end{wrapfigure} \noindent Joshua Berne serves at Bloomberg LP as a senior
software engineer on Bloomberg’s core library team. After the difficult choice
to pursue a career in software engineering over research mathematics, he has
been an active programmer in the financial industry, writing day trading
applications in C++ for E*TRADE Capital Markets and, after that, architecting
large distributed trading systems in Java for Instinet and IDC. Since joining
Bloomberg in 2017, he has been an active participant in the C++ Standards
Committee, seeking to bring the advancements made within Bloomberg to the
C++ Standard and thus to the rest of the world.  His first WG21 paper was
\cite{wg21P1332R0}

                  % Add about the author page


\cleardoublepage
%%%%%%%The front matter ends

%%%%%%%%%% Now we'll start on the main text for the book.
%% Main Text:
\pagenumbering{arabic}     % Arabic numbering
\chapter{Application Developers}

\section{What is an Allocator-Aware Type?}
\subsection{Defining a PMR Allocator-Aware Type}
\subsection{std::pmr Collections}

\section{Using Allocator-Aware Types}
\subsection{How to use a Custom Memory Resource}
\subsection{How to Choose an Allocator}
\subsection{Testing Code that Allocates}

\section{Case Study 1: Unique Value Counting}


../../src/books/cawp/cawp-library-writers.tex
../../src/books/cawp/cawp-writing-allocators.tex
\chapter{Advanced}
\section{Modern Hardware}
\section{Effective Benchmarking}

Here we would be discussing the approach we have to benchmarking.  \todo{Determine a better location for benchmarking section}

\begin{cppcodeblock}{awpadv_sample1_h}
  struct S {
    void foo();
  };
\end{cppcodeblock}

\begin{cppcodeblock}{awpadv_sample1_cpp}
  void S::foo()
  {
  }
\end{cppcodeblock}


% Repeat this command for each part, chapter, or glossary file.
%
%%%%% End main text


%%%%%%%%%% Back matter starts
%
% If your book contains appendices, you will need to add those here.

\cleardoublepage
\bibliographystyle{ieeetr-hacked}  % Your production editor and copy editor may want you to
                                                          % use a different bibliography style file.
                                                          % You can input that style file here.
\bibliography{cawp-bib}  % Loading the bibliography file.
\nocite{*}  % \nocite is used because no citations appear in this sample text

\cleardoublepage
\printindex   % If your book will include a LaTeX-generated index, you'll include it with this line.
%
../../src/books/cawp/cawp-appendix.tex
%%%%%%% Back matter ends


\end{document}

